\section{Experimental Setup}
\label{sec:experimental_setup}
\label{sec:setup}

\textbf{Models.} LLaMA-2-13B, Qwen3-8B and Gemma~4 E4B-it (4.5B effective parameters), each fine-tuned with LoRA ($r{=}16$, $\alpha{=}32$, 3 epochs) on 13{,}211 student-written PeerWise explanations \citep{bao2025exploring} pooled from five courses: Cardiff Biology, Sydney Biology, Auckland Law, and UK Medicine Years 1 and 2.

\textbf{Preference data.} For each base model we sample $n{=}3$ explanations for 500 training questions drawn from the same pooled corpus, so every pool contains questions from all five courses (Appendix~\ref{sec:pair_stats}); no pool shares a question with a test set. We keep the best- vs.\ worst-scoring pair when the score gap is $\geq 0.05$, which leaves 143--632 pairs. DPO: 5 epochs, lr $5{\times}10^{-5}$, $\beta{=}0.1$, effective batch 8.

\textbf{Test sets and decoding.} Each course has a fixed 100-question test set. Qwen3 and Gemma~4 decode greedily. The LLaMA-2 harness samples ($T{=}0.5$), so the same LLaMA-2 SFT checkpoint scores differently from run to run (e.g.\ Cardiff NLI 0.045--0.21 over 11 runs, Appendix~\ref{sec:provenance}). We therefore always compare a DPO policy with the SFT row of the \emph{same} evaluation run, and report a 95\% bootstrap confidence interval for the change. We have one DPO training seed per cell; the SimPO result of \S\ref{sec:reasoning-reward} has three.

\textbf{Metrics.} \emph{NLI}: $S_\text{NLI}$ with the training scorer. \emph{Held-out NLI}: the same with \texttt{cross-encoder/nli-deberta-v3-large}, which is used nowhere in training, pair construction or model selection. \emph{Copies}: the share of explanations that contain the correct-option text verbatim. ACR, BLEU, BERTScore and the verifier score are in Appendix~\ref{sec:per_domain_appendix}.

\textbf{Answer-masking test.}\label{sec:masking} An explanation can entail the answer by restating it. To see how much entailment remains without that, we replace, in the explanation, the full option text and every word beginning with one of the option's key terms (words of $\geq$4 letters) with \texttt{[MASK]}, and re-score with the held-out model. For example, with the option ``ATP is generated anaerobically through substrate level phosphorylation using creatine phosphate'', the explanation sentence ``ATP is generated anaerobically \ldots'' becomes ``[MASK] is [MASK] [MASK] \ldots'', while the reasons around it are kept. Masked scores are low for every model; what matters is whether DPO raises or lowers them relative to SFT.

\section{Results}
\label{sec:results}

\subsection{Entailment gains, and what they are made of}
\label{sec:main_results}

\begin{table}[t]
\centering
\caption{SFT$\to$RLearner-LLM on the five test sets. Each pair of numbers comes from \emph{one} evaluation run on the same 100 questions (LLaMA-2 decodes by sampling at $T{=}0.5$, so its SFT score varies between runs; Qwen3 and Gemma~4 decode greedily). \emph{NLI}: training scorer (DeBERTa-v3-small), bold if the 95\% bootstrap CI of the change excludes zero. \emph{Held-out}: DeBERTa-v3-large, used nowhere in training. \emph{Masked}: held-out NLI after every key term of the correct option is replaced by \texttt{[MASK]} in the explanation (\S\ref{sec:masking}); bold marks an increase. \emph{Copies}: \% of explanations containing the correct-option text verbatim.}
\label{tab:main}
\resizebox{\textwidth}{!}{%
\begin{tabular}{llcccccc}
\toprule
Model & Test set & NLI & 95\% CI of $\Delta$ & Held-out NLI & Masked held-out NLI & Words & Copies (\%) \\
\midrule
LLaMA-2-13B & Cardiff Bio. & \textbf{.082$\to$.319} & [+0.14, +0.33] & \textbf{.308$\to$.383} & .198$\to$.104 & 228$\to$80 & 54$\to$71 \\
LLaMA-2-13B & Sydney Bio. & \textbf{.069$\to$.356} & [+0.19, +0.38] & \textbf{.179$\to$.464} & .104$\to$.065 & 288$\to$55 & 37$\to$61 \\
LLaMA-2-13B & Auckland Law & .289$\to$.323 & [-0.07, +0.14] & .520$\to$.376 & .283$\to$.132 & 191$\to$80 & 46$\to$60 \\
LLaMA-2-13B & UK Med Y1 & \textbf{.069$\to$.425} & [+0.26, +0.45] & \textbf{.254$\to$.455} & .149$\to$.124 & 223$\to$69 & 57$\to$74 \\
LLaMA-2-13B & UK Med Y2 & \textbf{.076$\to$.388} & [+0.21, +0.41] & \textbf{.198$\to$.430} & .144$\to$.079 & 215$\to$74 & 49$\to$75 \\
\midrule
Qwen3-8B & Cardiff Bio. & .196$\to$.182 & [-0.11, +0.08] & .281$\to$.238 & \textbf{.071$\to$.141} & 59$\to$181 & 59$\to$78 \\
Qwen3-8B & Sydney Bio. & .174$\to$.231 & [-0.04, +0.15] & .320$\to$.290 & \textbf{.067$\to$.156} & 66$\to$179 & 40$\to$71 \\
Qwen3-8B & Auckland Law & .319$\to$.230 & [-0.20, +0.02] & \textbf{.297$\to$.409} & \textbf{.071$\to$.221} & 32$\to$199 & 41$\to$85 \\
Qwen3-8B & UK Med Y1 & .246$\to$.207 & [-0.13, +0.06] & .315$\to$.258 & \textbf{.081$\to$.162} & 56$\to$184 & 65$\to$75 \\
Qwen3-8B & UK Med Y2 & .163$\to$.207 & [-0.05, +0.13] & \textbf{.202$\to$.255} & \textbf{.073$\to$.151} & 74$\to$182 & 44$\to$75 \\
\midrule
Gemma~4 E4B & Cardiff Bio. & \textbf{.212$\to$.350} & [+0.03, +0.25] & \textbf{.239$\to$.400} & \textbf{.060$\to$.105} & 61$\to$45 & 57$\to$53 \\
Gemma~4 E4B & Sydney Bio. & .247$\to$.231 & [-0.12, +0.09] & .414$\to$.355 & \textbf{.081$\to$.110} & 61$\to$49 & 44$\to$38 \\
Gemma~4 E4B & Auckland Law & .391$\to$.438 & [-0.08, +0.17] & \textbf{.349$\to$.454} & \textbf{.057$\to$.083} & 29$\to$28 & 46$\to$48 \\
Gemma~4 E4B & UK Med Y1 & .296$\to$.391 & [-0.02, +0.21] & \textbf{.266$\to$.502} & \textbf{.042$\to$.106} & 55$\to$40 & 66$\to$51 \\
Gemma~4 E4B & UK Med Y2 & \textbf{.160$\to$.389} & [+0.12, +0.33] & \textbf{.264$\to$.412} & \textbf{.069$\to$.083} & 67$\to$54 & 49$\to$47 \\
\bottomrule
\end{tabular}}
\end{table}


Table~\ref{tab:main} gives the 15 (model, test set) cells; Figure~\ref{fig:masking} plots the held-out and masked changes. The three base models behave differently.

\textbf{LLaMA-2-13B: the gain is restatement.} Entailment rises in four of five test sets with confidence intervals well above zero (e.g.\ UK Med Y1 0.069$\to$0.425), and the held-out scorer agrees. But DPO also cuts the output from 190--290 to 55--80 words and raises the share that copies the answer (e.g.\ 49\%$\to$75\% on Med Y2). Once the answer is masked, the held-out score is \emph{lower} than SFT's in all five test sets. The model has learned to state the answer in a short output, not to argue for it. Auckland Law, the one non-significant cell, is also where ILearner-LLM, which refines each answer five times, stays ahead (Appendix~\ref{sec:per_domain_appendix}).

\textbf{Qwen3-8B: no reliable change.} No confidence interval excludes zero, the two scorers disagree on direction in two cells, and outputs triple in length (32--74$\to$179--199 words). The masked score rises, but with three times as much text to mask around we do not read this as evidence either way.

\textbf{Gemma~4 E4B: a gain that survives masking.} The held-out scorer improves in four of five test sets (all but Sydney) and the masked score in all five, while outputs get \emph{shorter} (by 1--16 words) and copy the answer no more often (change $-15$ to $+2$ points). The training-scorer change has a CI above zero on Cardiff and Med Y2; on Med Y1 and Law the CI touches zero. This is the result we consider robust: on Gemma~4, the entailment reward improves how well the explanation supports the answer beyond repeating it. Whether reasoning judges also prefer these outputs is a separate question, which \S\ref{sec:single_dual} answers in the negative. On the strict question-quality filter of Appendix~\ref{sec:tierc_appendix} the Gemma~4 gains are larger (held-out 0.32$\to$0.56 on Cardiff, 0.37$\to$0.55 on Sydney).

Across all 15 cells, the CI excludes zero in 6, a one-sided paired Wilcoxon test gives $p{<}0.05$ in 9, and the two scorers agree on the direction of the change in 12.

\subsection{Does combining NLI with a verifier help?}
\label{sec:single_dual}
\label{sec:variant_ablation}

\begin{table}[t]
\centering
\caption{Gemma~4: every reward variant trained on pairs drawn from \emph{one} shared pool (500 prompts $\times$ 3 candidates, each scored once by every signal), with identical DPO settings. Means over the five test sets (500 questions, greedy decoding). Rows with $\pm$ give mean$\pm$s.d.\ over three DPO seeds; all other cells are seed 42. The CI is a paired bootstrap of the seed-42 change in masked held-out NLI vs.\ SFT. Ver: \emph{reproduced} verifier (Appendix~\ref{sec:verifier_repro}). Judge: \% of non-tie pairwise reasoning judgements won against SFT (Qwen2.5-7B / Llama-3.1-8B / Gemma-4-31B, 500 items each, order-flipped). Round-2 SimPO rows use $\beta{=}2$, $\gamma{=}1$, 3 epochs; the reasoning-judge reward is a fourth model (Qwen3-8B) that is never used for evaluation. Judge cells of three-seed rows pool the three seeds' judgements.}
\label{tab:g4shared}
\resizebox{\textwidth}{!}{%
\begin{tabular}{lccccccccc}
\toprule
Reward & Pairs & NLI & Held-out NLI & Masked held-out & 95\% CI $\Delta$masked & Words & Copies (\%) & Ver & Judge vs SFT \\
\midrule
SFT & -- & 0.261 & 0.307 & 0.062 & -- & 55 & 52 & 2.89 & -- \\
\midrule
Verifier only & 314 & 0.160 & 0.266 & 0.086 & [+0.001, +0.048] & 99 & 51 & 3.18 & -- \\
$H_A$, $w{=}0.25$ & 258 & 0.287 & 0.397 & 0.087 & [+0.003, +0.047] & 66 & 54 & 3.00 & -- \\
$H_A$, $w{=}0.5$ & 257 & 0.335 & 0.435 & 0.103 & [+0.019, +0.064] & 62 & 52 & 2.95 & -- \\
$H_A$, $w{=}0.75$ & 262 & 0.331 & 0.425 & 0.100 & [+0.016, +0.060] & 57 & 51 & 2.92 & -- \\
NLI only & 270 & 0.381 & 0.473 & 0.133 & [+0.048, +0.095] & 53 & 47 & 2.89 & 39 / 42 / 41 \\
$H_M$ (Table~\ref{tab:main} reward) & 159 & 0.311$\pm$0.011 & 0.405$\pm$0.015 & 0.086$\pm$0.006 & [+0.008, +0.047] & 49 & 50 & 2.89 & 38 / 41 / 39 \\
$H_M$, dominating pairs only & 117 & 0.270 & 0.334 & 0.080 & [-0.002, +0.038] & 60 & 53 & 2.94 & -- \\
Masked NLI & 93 & 0.272$\pm$0.006 & 0.343$\pm$0.011 & 0.084$\pm$0.007 & [-0.002, +0.036] & 58 & 54 & 2.91 & 47 / 47 / 47 \\
NLI only, 93 random pairs & 93 & 0.277$\pm$0.010 & 0.361$\pm$0.004 & 0.083$\pm$0.010 & [-0.004, +0.033] & 56 & 54 & 2.88 & 46 / 53 / 53 \\
\midrule
\multicolumn{10}{l}{\emph{Round 2: lighter NLI training, SimPO, and a reasoning-judge reward (Qwen3-8B) on the same pool}} \\
NLI only, 1 epoch & 270 & 0.240 & 0.300 & 0.077 & [+0.004, +0.028] & 54 & 55 & 2.89 & 50 / 49 / 52 \\
NLI only, 2 epochs & 270 & 0.310 & 0.351 & 0.077 & [+0.001, +0.029] & 50 & 53 & 2.89 & 49 / 48 / 47 \\
NLI only, $\beta{=}0.5$ & 270 & 0.287 & 0.339 & 0.085 & [+0.005, +0.041] & 53 & 54 & 2.90 & 48 / 49 / 50 \\
NLI only (SimPO) & 270 & 0.391 & 0.403 & 0.092 & [+0.009, +0.053] & 29 & 44 & 2.84 & 39 / 36 / 31 \\
Reasoning judge (SimPO) & 293 & 0.244 & 0.313 & 0.084 & [+0.001, +0.045] & 60 & 45 & 2.90 & 48 / 60 / 61 \\
Judge + NLI (SimPO) & 390 & 0.343 & 0.410 & 0.115 & [+0.028, +0.078] & 48 & 47 & 2.86 & 42 / 47 / 40 \\
\midrule
\multicolumn{10}{l}{\emph{Round 3: judge pairs restricted to candidates within 15 words of each other}} \\
Judge, length-matched pairs (SimPO) & 182 & 0.258$\pm$0.011 & 0.316$\pm$0.000 & 0.075$\pm$0.001 & [-0.006, +0.030] & 54 & 52 & 2.88 & 52 / 56 / 54 \\
Judge + NLI, length-matched (SimPO) & 254 & 0.288 & 0.364 & 0.096 & [+0.011, +0.057] & 53 & 52 & 2.88 & 46 / 52 / 50 \\
\bottomrule
\end{tabular}}
\end{table}


Our original motivation for adding the verifier was that NLI alone would collapse into short, incoherent text. We did not observe this. With LLaMA-2, NLI-only DPO produces 71--97-word outputs with the same verifier scores as the combined reward, and the same entailment gains (+0.24 vs.\ +0.29 on Cardiff, +0.29 vs.\ +0.29 on Sydney, +0.31 vs.\ +0.31 on Med Y2); verifier-only DPO, with no NLI signal at all, gives +0.24 and +0.22 (Appendix~\ref{sec:single_dual_appendix}, Table~\ref{tab:single_dual}). On Qwen3 the combined, NLI-only and verifier-only policies are indistinguishable. Those runs used separately sampled candidate pools. To compare the rewards cleanly we built a single pool of scored candidates for Gemma~4 and trained every variant of Table~\ref{tab:rewards} on pairs drawn from it (Table~\ref{tab:g4shared}).

\GSHAREDTEXT

\subsection{Rewarding reasoning directly}
\label{sec:reasoning-reward}

Entailment, masked or not, does not ask whether the reason given is valid. We therefore also train with a reward that does. From the Cardiff Tier-C training prompts we sample three explanations per question from the Gemma~4 SFT model and rank them with a pairwise judge (DeepSeek-R1-32B) asked which explanation better justifies \emph{why} the stated answer is correct, ignoring length and formatting, with randomised order. Of 321 pairs we keep the 112 whose lengths differ by at most 15 words (chosen 46, rejected 44 words), so length cannot drive the preference; a judge from another family (Llama-3-8B) agrees on 69\% of them. We train with SimPO, whose length-normalised objective removes the incentive to lengthen; the policy writes 68 words on the test set against 66 for SFT.

\begin{table}[t]
\centering
\caption{Reasoning reward vs.\ NLI reward on the 191 Cardiff Tier-C test questions (Gemma~4, same SFT checkpoint, greedy decoding). \textbf{Top:} pairwise reasoning judgements (first policy wins / second wins / ties; \% of non-tied pairs; two-sided sign test), same prompt and order flip as Table~\ref{tab:reasoning_judges}, judges run locally in bf16. \textbf{Bottom:} the same outputs on the entailment metrics.}
\label{tab:reasoning_vs_hybrid}
\resizebox{\textwidth}{!}{%
\begin{tabular}{lccc}
\toprule
Judge & SimPO vs SFT & Hybrid-DPO vs SFT & SimPO vs Hybrid-DPO \\
\midrule
Qwen2.5-7B-Instruct & 138/49/4 (74\%, $p{=}5{\times}10^{-11}$) & 96/94/0 (51\%, $p{=}0.94$) & 134/57/0 (70\%, $p{=}2{\times}10^{-8}$) \\
Llama-3.1-8B-Instruct & 160/30/0 (84\%, $p{=}1{\times}10^{-22}$) & 111/79/0 (58\%, $p{=}0.02$) & 148/43/0 (77\%, $p{=}1{\times}10^{-14}$) \\
\bottomrule
\end{tabular}}

\vspace{4pt}
\resizebox{0.86\textwidth}{!}{%
\begin{tabular}{lccccc}
\toprule
Policy & Words & Copies (\%) & NLI (train) & NLI (held-out) & Masked held-out NLI \\
\midrule
SFT & 66 & 53 & 0.186 & 0.321 & 0.086 \\
Hybrid-DPO (NLI reward) & 78 & 71 & 0.390 & 0.558 & 0.136 \\
SimPO (reasoning reward) & 68 & 51 & 0.165 & 0.281 & 0.049 \\
DPO (reasoning reward) & 60 & 51 & 0.202 & 0.289 & 0.073 \\
Self-distillation (reasoning reward) & 44 & 58 & 0.229 & 0.343 & 0.064 \\
\bottomrule
\end{tabular}}
\end{table}


On all 191 Cardiff Tier-C test questions, three judges independent of the reward (Qwen2.5-7B, Llama-3-8B and gpt-5-mini) prefer SimPO's explanations to SFT's in 70\%, 72\% and 82\% of non-tied pairs ($p\ll10^{-3}$, Appendix~\ref{sec:reasoning_appendix}); three training seeds give $76.4\pm0.8\%$ under Qwen2.5. In a blinded study, four raters with biology or medical background prefer SimPO in 63\% of non-tied judgements on 24 length-matched pairs ($p\approx0.02$; agreement between raters is modest, 49\%; Appendix~\ref{sec:human_appendix}).

Table~\ref{tab:reasoning_vs_hybrid} puts the NLI-reward policy, trained from the same SFT model on the same course, into this comparison. The two rewards do different things. The NLI-reward policy doubles entailment (0.19$\to$0.39) and part of that survives masking, but the reasoning judges prefer it to SFT only 51\% and 58\% of the time, and prefer SimPO to it 70\% and 77\% of the time. The reasoning reward leaves entailment unchanged or lower. Higher entailment is not better-judged reasoning; each reward improves what it measures.

\subsection{LLM judges prefer the longer output}
\label{sec:verbosity}
We also asked GPT-4o-mini to compare explanations pairwise (one call per item, order randomised; Appendix~\ref{sec:pairwise_appendix}). In every comparison it chose the longer output: LLaMA-2 SFT (234 words) over verifier-only DPO (88) in 69\% of items, Qwen3 combined-reward DPO (183) over its SFT (59) in 95\%, and its own outputs (237) over Qwen3 DPO (183) in 95\%. Because length differs in each of these comparisons, they say nothing about quality; we report them as a reason to use length-matched comparisons and several judges, as in \S\ref{sec:reasoning-reward}.
